\documentclass[11pt]{article}
\usepackage[T1]{fontenc}
\usepackage{lmodern,amsmath,amssymb,amsthm,booktabs,microtype}
\usepackage[margin=0.9in]{geometry}
\usepackage[hidelinks]{hyperref}
\newtheorem{lemma}{Lemma}
\newtheorem{theorem}{Theorem}
\newtheorem{corollary}{Corollary}
\newcommand{\B}{\mathsf B}
\newcommand{\Con}{\operatorname{Con}}
\title{Two contextual state folds and a 246-state blank-start machine}
\author{ProveIt research contribution}
\date{17 October 2026}
\begin{document}
\maketitle

\begin{abstract}
The prefix250 register interpreter has two states whose recurrent execution
uses only one read-symbol branch. Their live branches coincide with retained
states. Folding these two states also disconnects two states used by its old
bootstrap. The resulting recurrent interpreter has 246 working states. Its
blank initialization needs no extra state: starting in the existing carry
state 95 reaches a safe root after exactly 410,735 literal transitions.
This note proves the state-folding transport. It retains the mathematical
and executable-certificate proof status of prefix250; it does not claim a
new proof-assistant endpoint.
\end{abstract}

\section{Exact construction}

All state numbers in the proof are the original labels, not canonical
renumberings. The baseline is the exact file
\texttt{zfc\_prefix250\_old.tm}, SHA-256
\begin{quote}\small\ttfamily
fd455304bad82a6da7ab83149352d0b1db27557582c4cc6312de003c96ff4e43.
\end{quote}
Delete states $0,1,3,97$. Change only the following two surviving rows:
\begin{center}
\begin{tabular}{rrrrl}
\toprule
State & Read & Write & Direction & Next state\\
\midrule
60 & 0 & 0 & R & 20\\
255 & 0 & 1 & L & 255\\
\bottomrule
\end{tabular}
\end{center}
Every other surviving transition is unchanged. Select state 95 as the blank
start; state 254 remains the separate halt label. The output has 246 working
states, 492 binary transitions, and no edge to a deleted state. Its
original-label SHA-256 is
\begin{quote}\small\ttfamily
303911c554284c49ce09042d8127feba4da260baae766e32260034aa3bdb82ee.
\end{quote}
The canonical table gives state 95 number zero, assigns the other surviving
labels in increasing order, and gives the halt label number 246. Its
SHA-256 is
\begin{quote}\small\ttfamily
89cda615b3ea9fcac07ad6b08e0351faa34841a56935b6f5846d3920f13bc0fc.
\end{quote}
All 246 working states are reachable in the directed transition graph.
This graph fact is distinct from reachability on the actual blank run.

\section{The one-step quotient identity}

Write $T$ for prefix250 and $U$ for the new table. A transition
$(w,d,q')$ writes the symbol $w$, moves exactly one tape cell in direction
$d$, and enters $q'$. Define the state map
\[
 f(0)=20,\qquad f(97)=255,\qquad f(q)=q\quad(q\notin\{0,97\}).
\]
The excluded state/read pairs are
\[
 E=\{(0,0),(1,0),(1,1),(3,0),(3,1),(97,0)\}.
\]
The domain restriction is essential. In particular the old blank start
reads zero in state 97, so that old bootstrap cannot be transported by
this quotient.

\begin{lemma}[Literal commutation]\label{lem:commute}
For every transition $(q,b)\notin E$ of $T$, if
$T(q,b)=(w,d,q')$, then
\[
 U(f(q),b)=(w,d,f(q')).
\]
\end{lemma}
\begin{proof}
The only two incoming-target changes are
$60/0:0R0\mapsto0R20$ and $255/0:1L97\mapsto1L255$.
The surviving live rows of the folded states agree exactly with their
representatives:
\[
 T(0,1)=T(20,1)=(0,L,2),\qquad
 T(97,1)=T(255,1)=(1,L,131).
\]
Every other admitted row is unchanged, including both rows of states 20
and 255 except for the explicitly retargeted 255/0 row. Direct inspection
of the complete table excludes any other edge to 0 or 97, and excludes
edges to 1 or 3 outside the pairs in $E$. These observations cover all
494 admitted original state/read pairs. The regeneration script also
checks the displayed identity separately on each of those 494 pairs.
\end{proof}

\begin{corollary}[Exact-time trace transport]\label{cor:trace}
Let a finite $T$ execution avoid $E$. There is a $U$ execution of the same
length whose tape and head coordinate agree at every time and whose state
at every time is the image under $f$ of the original state. In particular,
the original and folded executions have the same halting outcome on such
a trace, including the final write and head movement.
\end{corollary}
\begin{proof}
Induct on the number of transitions, using Lemma~\ref{lem:commute} at
each step. The written symbol and movement are identical, so the next
read symbol also agrees. The halt label is fixed by $f$.
\end{proof}

\section{Why every recurrent excursion avoids the excluded pairs}

The physical register file for a finite list $v$ is
\[
 \B(v)=\prod_{i<|v|} 0\,1^{v_i+1}.
\]
Each block therefore contains at least one one-cell, even when its
numerical value is zero. The runtime invariant has nine primary blocks
and at least one spare block. Register values and the finite spare suffix
are otherwise arbitrary. The original local identities admit arbitrary
left tape context.

\begin{lemma}[Marking fold]\label{lem:mark}
The selector/marking excursion on any admitted unmarked register file
enters old state 0 only while reading one. It never enters states 1 or 3.
Replacing its final selector edge by $60/0:0R20$ preserves the excursion's
whole tape, head coordinate, exit state, and number of transitions.
\end{lemma}
\begin{proof}
The selector chain passes exactly the requested number of ordinary blocks.
Its final state 60 reads the zero separator of the selected block and
moves right to state 0. The next cell is the compulsory first one of
that block; this is true uniformly for every natural register value.
Thus state 0 takes its read-one row, erases that cell to mark the register,
and enters state 2 one cell to the left. The corresponding read-one row
of state 20 is identical. Both executions then have exactly the same
configuration, not merely equivalent register contents.

The only outside entry into the pair of states $\{1,3\}$ is old state
0's read-zero branch. That branch was just excluded. Neither state 97
nor the carry termination circuit occurs in this marking excursion.
\end{proof}

\begin{lemma}[Pulse fold]\label{lem:pulse}
Each increment, positive-decrement, and zero-decrement excursion from
its admitted marked register file enters old state 97 only while reading
one. Replacing $255/0:1L97$ by $255/0:1L255$ preserves every such excursion's
whole tape, head coordinate, exit state, and exact duration.
\end{lemma}
\begin{proof}
In these excursions, the only entry into old state 97 is the read-zero
branch of state 255. State 255 in turn is entered from state 75's
read-zero branch. Let $c$ be the head coordinate at that state-75 step.
The actual rows are
\[
 T(75,0)=(1,R,255),\qquad T(255,0)=(1,L,97).
\]
The first step writes one at $c$ and moves to $c+1$. If the second row
is used, it writes at $c+1$ and returns to $c$. Nothing has changed the
one at $c$. Consequently old state 97 necessarily reads one. Its row
is $(1,L,131)$, exactly the read-one row of state 255. The substituted
three-step sequence therefore has the same endpoint and the same tape:
\[
 75\to255\to97\to131
 \quad\longmapsto\quad
 75\to255\to255\to131.
\]
The read-one branch from the first visit to state 255 is unchanged.
The earlier rightward carries, reverse carries, seekers, and return scans
are unchanged. These excursions never use states 0, 1, or 3.
\end{proof}

The old exact local identities therefore transport without any correction
to their register effects, boundary shifts, or timing formulas. Increment
still appends a spare register of value one. Positive decrement still
shifts the file boundary one cell right. Zero decrement still restores
the original file and exits in state 112 at its old left boundary. The
nonempty spare invariant is preserved.

Outside these register macros, the complete existing 8,391 instruction
guard cases execute none of the pairs in $E$. This statement can be
checked by recording literal transition lookups in the supplied affine
symbolic interpreter. The check is an overapproximation: lookups used to
test whether a whole run is a self-loop may be included as well. None
of the excluded pairs is permitted by that larger lookup set. Applying
the transported local contracts gives exactly the same complete
instruction targets, for every affine gap parameter and every register
vector satisfying its named guard. Thus the table folds do not change
the source program, register permutation, PC transition, or logical
calibration.

\section{New blank bootstrap}

\begin{lemma}[Complete safe root]\label{lem:boot}
Starting on the all-zero bi-infinite tape in old state 95 with head at
coordinate zero, the new machine reaches old state 96 after exactly
410,735 transitions, without an earlier halt. The head is at $-5$,
the marker is at $-6$, the header is at $176$, and the PC is empty.
The complete tape is
\[
 0^\infty\,1\,0^{181}\,1\,
 \B([0,0,0,3,0,0,3,0,0]\mathbin{+\!+}[1]^{167})\,0^\infty.
\]
There are exactly 351 one-cells; every other cell is zero.
\end{lemma}
\begin{proof}
This is an exact finite computation reproduced by \texttt{build246.py}.
The implementation starts with an empty set of one-cell coordinates,
applies each literal row, and checks both the final state/head pair and
the entire resulting set against the displayed tape. No prespecified
window or omitted exterior cells occur in this comparison. An independent
array-tape search and a separate set-tape replay agree on this result.
The evidence file lists every one-cell coordinate.
\end{proof}

Under the unchanged logical-to-physical permutation
$(6,2,0,4,3,7,8,5,1)$, the nine physical values correspond to
\[
 (N,T,W,A,D,H,X,Y,L)=(3,0,0,0,3,0,0,0,0).
\]
The three formula fields $T,W,H$ are all zero; the other six fields are
permitted arbitrary natural values by the existing safe dirty-root
theorem. In particular the nonzero values of $N$ and $D$ do not require
a new frontend argument. The header span is $181\ge40$, and there are
167 spare blocks, so the complete recurrent invariant holds.

\begin{theorem}[Runtime and consistency transport]
Assume the prefix250 local-instruction and safe-source-root calibration
proved in its accompanying report and the unchanged ProveIt sources.
The displayed 246-state blank-start machine halts if and only if the
unchanged source program halts at the safe root
$(3,0,0,0,3,0,0,0,0)$. Consequently the same ordinary mathematical
calibration gives
\[
 \mathrm{ZFC}\ \vdash\
 \bigl(\text{the 246-state machine does not halt from blank tape}\bigr)
 \longleftrightarrow \Con(\mathrm{ZFC}).
\]
\end{theorem}
\begin{proof}
Lemma~\ref{lem:boot} establishes a complete safe root. The two local
fold lemmas and the admitted instruction traces preserve the entire
instruction simulation. Each segment has exactly its original positive
number of literal steps, so the original no-spurious-halt argument
applies in both directions. Induction on successive source instructions
preserves the span and nonempty-spare invariant. The safe-root theorem
then identifies source halting with native-theory refutation, and the
unchanged set-theoretic calibration identifies that refutation with
ZFC inconsistency.
\end{proof}

The dependency in this theorem is stated explicitly. The contribution
is the new 246-state table, the universally valid contextual quotient,
and a checked new blank bootstrap. It does not replace the prior
calibration, and it does not infer object-theoretic ZFC provability from
a proof assistant's ambient metatheory or a dependency report.

\end{document}
